\subsection{The bordered landing obstruction and its exact repair}
\label{sec:A1-bordered}

There is a separate issue before the size of the tensor range is
considered. Cote 67/065--066 defines admissibility by comparing the
circle actions of the discs with those of the hyperbolic metric on
their complement. The proposed target in 67/075--076 retains an
admissible system only at the external markings. On a nodal curve,
the direct extension of that definition is componentwise on the
normalization, with the branches of the nodes left as punctures.
The following obstruction uses precisely this interpretation.

\begin{proposition}[Refuted: the bordered map lands in the stated target]
\label{A1-bordered-landing}
Canonical cutting and capping need not produce an admissible system
of external discs on the normalization when the node caps are
forgotten and their centers are left as punctures. This already
fails for a four-holed sphere with one separating cut, all five
geodesic lengths---the four external boundaries and the cut---equal
to $2\pi$. Thus, with the componentwise admissibility just specified,
the displayed arrow $MDB\longrightarrow M_{\mathrm{st}}ND$ of
67/076 is not a well-defined map on all objects.
\end{proposition}
\begin{proof}
Take two hyperbolic pants whose three boundary lengths are $2\pi$,
and glue one boundary of each by any orientation-reversing arclength
isometry. The result is a hyperbolic four-holed sphere with the
indicated separating geodesic. Cutting and canonically capping gives
two marked spheres. On either one, label the external centers $0,1$
and the center of the node cap $\infty$. By
Proposition~\ref{A2-schwarzian}, its projective-connection
coefficients are $a_0=a_1=a_\infty=1$.

Let $u_0$ and $u_1$ be centered unit-disc coordinates on the two
external caps, and use the local coordinates $z$ and $\zeta=1-z$,
respectively. Write
\[
 u_0(z)=c_0z(1+b_0z+O(z^2)),\qquad
 u_1(1-\zeta)=c_1\zeta(1+b_1\zeta+O(\zeta^2)).
\]
The formula in that proposition, first as written and then with the
labels $0,1$ interchanged, gives $b_0=b_1=1/2$. These coefficients
depend only on the two embedded pointed discs: changing either
centered unit-disc coordinate multiplies it by a unit scalar and
does not change $b_i$.

Suppose that after forgetting the cap at $\infty$, the two retained
discs were admissible on the twice-bordered, once-punctured sphere.
Its hyperbolic metric would have a cusp at $\infty$ and positive
boundary lengths $L'_0,L'_1$ at the retained discs. Set
\[
 a'_0=\frac12\left(1+\frac{(L'_0)^2}{4\pi^2}\right)>\frac12,
 \qquad
 a'_1=\frac12\left(1+\frac{(L'_1)^2}{4\pi^2}\right)>\frac12,
 \qquad a'_\infty=\frac12.
\]
The same Schwarzian formula applies when a boundary is a cusp: a
cusp developing coordinate is a constant multiple of $\log q$,
whose Schwarzian is $1/(2q^2)$. Admissibility would identify the
canonical cap domains with the two already given discs. Their
unchanged coefficients would therefore satisfy
\[
 \frac{a'_0+a'_1-1/2}{2a'_0}=\frac12,
 \qquad
 \frac{a'_0+a'_1-1/2}{2a'_1}=\frac12.
\]
The first equality gives $a'_1=1/2$ and the second $a'_0=1/2$,
contradicting both positive boundary lengths. This is an obstruction
to the full disc shapes; it allows the new lengths to vary freely.
\end{proof}

The word ``admissible'' could instead be assigned a new meaning on
a nodal curve, namely joint extendibility by suitable node caps.
That is a possible repair, but it is a different condition from
admissibility of just the external system on the punctured
normalization. The next theorem describes the repaired data and
their image without retaining the node discs as independent
parameters.

Let $L$ be the set of actual external boundary components; other
external markings, if present, remain cusps. For each vertex $v$,
let $J_v$ contain all node branches and all legs in $L$ incident to
$v$. Put $P=E\sqcup L$ and define
\[
 B_v:\R^P\longrightarrow\R^{J_v},\qquad
 (B_vt)_a=
 \begin{cases}
 t_e,&a\text{ is a branch of }e,\\
 t_l,&a\text{ is the external leg }l.
 \end{cases}
\]
In particular a loop edge contributes two entries, while a leg
contributes one. For the fixed marked normalization $Y=(Y_v)$,
let $S_v,K_v$ be the cap potential and support function with active
set $J_v$ and all remaining markings as cusps, and define
\[
 \Psi_{\Gamma,L}(t)=\sum_v S_v(B_vt),\qquad
 K_{\Gamma,L}(p)=\sum_v K_v(B_vp).
\]
An empty active set contributes zero. Choose tangent frames only
to write radii as positive numbers; the conditions below are
intrinsic under changes of those frames.

\begin{theorem}[Proved: exact bordered model by tensors and shadows]
\label{A1-bordered-shadow-range}
On the fixed graph and normalization, record the nonzero node
tensors $q_e$ and the external cap shadows, of radii $r_l>0$.
Canonical cutting and capping gives a real-analytic diffeomorphism
from bordered cut surfaces, with actual toric seam phases, onto
\begin{align*}
 \mathcal V_{\Gamma,L}(Y)=\biggl\{(q,r):\quad&
 \sum_{e\in E}p_e\log|q_e|+
 \sum_{l\in L}p_l\log r_l<K_{\Gamma,L}(p)\\
 &\text{for every }0\ne p\in[0,\infty)^P\biggr\}.
\end{align*}
The logarithmic radial image is open, convex, and downward closed.
The source length at every edge or external leg is
$2\pi\sqrt{t_j}$, where $t$ is the unique solution
\[
 \nabla\Psi_{\Gamma,L}(t)
   =\bigl((\log|q_e|)_{e\in E},(\log r_l)_{l\in L}\bigr).
\]
Thus the node tensors and the external shadows determine all node
and external cap domains, all boundary lengths, and the cut surface.
The statement includes self-nodes and fixed cusp markings. Its
tangent map is an isomorphism; as the normalization varies, the
range forms an open real-analytic subbundle. The inclusion into
the full nonzero-tensor and positive-shadow bundle is an
equivariant homotopy equivalence.
\end{theorem}
\begin{proof}
The chain rule and the seam tensor formula give
\[
 \frac{\partial\Psi_{\Gamma,L}}{\partial t_e}
       =\log r_{a_e}+\log r_{b_e}=\log|q_e|,
 \qquad
 \frac{\partial\Psi_{\Gamma,L}}{\partial t_l}=\log r_l.
\]
By Theorem~\ref{A2b-energy} its Hessian is
\[
 d^2\Psi_{\Gamma,L}=
       \sum_v B_v^T(d^2S_v)B_v>0.
\]
Indeed every coordinate of $\R^P$ occurs in at least one $B_v$,
so their combined map is injective. Integrating this positive
quadratic form along a segment proves global injectivity, and
the inverse function theorem proves local analytic invertibility.

Lemma~\ref{A2b-recession} gives recession function
$K_{\Gamma,L}$. The strict inequalities are necessary by the
componentwise inequalities of Theorem~\ref{A2b-image}. Conversely
let $x$ satisfy the displayed support inequalities. The convex
function $\Psi_{\Gamma,L}-x\cdot t$ has bounded sublevels: an
unbounded sublevel sequence, normalized in direction, would
produce $p\geq0$, $p\ne0$, with recession slope
$K_{\Gamma,L}(p)-p\cdot x\leq0$. Its lower semicontinuous closure
has a minimizer in the closed orthant. No finite minimizer lies
on a boundary face, because the inward diagonal derivative is a
sum of cap log radii with positive multiplicities; a vanishing
coordinate supplies at least one term tending uniformly to
$-\infty$ by the collar estimate, and all remaining terms have
uniform upper bounds. The unique minimizer is therefore interior
and satisfies $\nabla\Psi_{\Gamma,L}=x$. This proves the exact
radial image and its analytic inverse.

For those recovered parameters the canonical cap systems exist
uniquely by Theorem~\ref{A4-grafting}. Edge branches have equal
lengths by construction. The phase of $q_e$ determines their toric
identification, and Proposition~\ref{A1-retained-discs} reconstructs
the surface. The same actual-phase convention as in
Theorem~\ref{A1-graph-range} applies; fully marked real twists
require the corresponding covering on the target.

With varying normalization, the derivative is block triangular
with the identity on normalization parameters, the positive
Hessian on lengths, and the identity on phases. Hence it is an
isomorphism and the inverse varies real-analytically. Finally set
every $t_j=1$. The resulting canonical shadows define a positive
reference radius at each leg and a reference circle in every
tensor line. Linear interpolation of logarithmic radii to these
references is a deformation retraction of both the full radial
bundle and the stated convex subbundle onto the same phase-torus
bundle. It is intrinsic and respects automorphisms, as in
Theorem~\ref{A1-graph-range}. This proves the last assertions.
\end{proof}

\begin{corollary}[Proved: the image when full external discs are retained]
\label{A1-bordered-shape-image}
Replace the manuscript's external-only admissibility requirement
by arbitrary embedded pointed external discs with disjoint
analytic closures avoiding every other special point. The exact
image of canonical surgery in this target consists of the tuples
$(Y,(q_e),(D_l))$ satisfying both of the following conditions:
\begin{enumerate}
\item The tensors and the shadows $r_l=r(D_l)$ satisfy all support
inequalities of Theorem~\ref{A1-bordered-shadow-range}.
\item Let $t$ be the unique gradient inverse furnished by that
theorem. On every component, the given external discs $D_l$ are
exactly the external members of the canonical full cap system
with squared length parameters $B_vt$.
\end{enumerate}
These conditions determine a real-analytic graph over
$\mathcal V_{\Gamma,L}(Y)$. They give a complete reconstruction
criterion, and do not require the retained external subsystem
to be admissible after the node caps have been discarded.
\end{corollary}
\begin{proof}
Both conditions hold for a surgery output. Conversely condition
(1) uniquely constructs the full canonical cap systems, equal
paired lengths, and phase identifications. They weld by
Proposition~\ref{A1-retained-discs}; condition (2) says exactly
that their external members are the prescribed discs. The
constructed surface is therefore a preimage, and uniqueness
follows from Theorem~\ref{A1-bordered-shadow-range}. Analytic
dependence of the reconstructed disc embeddings is
Lemma~\ref{A2b-regularity}. Proposition~\ref{A1-bordered-landing}
shows why external-only admissibility cannot be imposed in
addition as an automatic property.
\end{proof}
